\chapter*{ESPECIFICACIONES DE MONTAJE}
\addcontentsline{toc}{chapter}{ESPECIFICACIONES DE MONTAJE}

% Reinicio de contadores para numeracion coherente con el capitulo 4.
\setcounter{section}{0}
\renewcommand{\thesection}{4.\arabic{section}.}
\renewcommand{\thesubsection}{4.\arabic{section}.\arabic{subsection}.}
\renewcommand{\thesubsubsection}{4.\arabic{section}.\arabic{subsection}.\arabic{subsubsection}.}
\renewcommand{\theHsection}{montaje.\arabic{section}}
\renewcommand{\theHsubsection}{montaje.\arabic{section}.\arabic{subsection}}
\renewcommand{\theHsubsubsection}{montaje.\arabic{section}.\arabic{subsection}.\arabic{subsubsection}}

\section{INSTALACION DE CANALIZACIONES}

\subsection{TUBERIAS EMPOTRADAS EN MUROS Y COLUMNAS}
Para la instalacion de tuberias en muros de ladrillo se deberan seguir las siguientes pautas constructivas:
\begin{itemize}
  \item \textbf{Canaleteado:} Las canaletas se ejecutaran utilizando amoladora con disco de corte para ladrillo y martillo/cincel, prohibiendose el uso de herramientas pesadas que comprometan la estabilidad estructural del muro.
  \item \textbf{Profundidad:} Las tuberias deben quedar recubiertas con una capa minima de \(15\text{ mm}\) de mortero o tarrajeo para evitar fisuras.
  \item \textbf{Curvas:} Se emplearan curvas de fabrica de \(90^\circ\) unidas con pegamento PVC, limitando el numero a un maximo de 3 curvas de \(90^\circ\) entre caja y caja.
\end{itemize}

\subsection{TUBERIAS EN LOSAS DE CONCRETO}
\begin{itemize}
  \item Las tuberias se fijaran firmemente sobre la malla de acero de refuerzo antes del vaciado del concreto, empleando alambre galvanizado.
  \item Se evitaran cruces de tuberias en zonas de alta concentracion de esfuerzos como encuentros de vigas y columnas.
  \item Todas las uniones seran verificadas para asegurar hermeticidad frente a la lechada de cemento durante el vaciado.
\end{itemize}

\section{INSTALACION DE CAJAS Y ALTURAS REGLAMENTARIAS}

Las cajas de salida e interruptores deben quedar perfectamente niveladas y enrasadas con el tarrajeo o acabado definitivo del muro (tolerancia maxima de embutido de \(6\text{ mm}\)). Se definen las siguientes alturas de montaje medidas desde el nivel de piso terminado (N.P.T.) al eje de la caja:

\begin{center}
\begin{tabular}{l c l}
\toprule
\textbf{Dispositivo / Salida} & \textbf{Altura (m)} & \textbf{Observaciones} \\
\midrule
Tablero General (TG) & 1.50 m & Medido al borde superior del gabinete \\
Tablero de Distribucion T2 & 1.50 m & En segundo piso \\
Interruptores de pared & 1.20 m & Cerca de las puertas, lado opuesto a bisagras \\
Tomacorrientes generales & 0.30 m & Bipolar doble con linea de tierra \\
Tomacorrientes sobre mesa & 1.10 m & En cocina amplia y zona de servicio \\
Tomacorrientes de bano & 1.10 m & Protegido contra salpicaduras \\
Cajas de paso & 1.80 m o techo & Ubicaciones accesibles para registro \\
Centros de alumbrado & Techo / Losa & Centrados en los ambientes \\
\bottomrule
\end{tabular}
\end{center}

\section{CABLEADO Y CONEXIONADO}

El paso de los conductores por el interior de las tuberias se realizara una vez concluido el tarrajeo de muros, vaciado de techos y cuando la obra civil este completamente seca.
\begin{itemize}
  \item \textbf{Guia de traccion:} Se utilizara cinta guia (wincha de acero o nylon) introduciendo los cables en un solo sentido sin realizar esfuerzos excesivos que puedan danar el aislamiento.
  \item \textbf{Lubricacion:} Se utilizara exclusivamente talco industrial o estearina como lubricante de arrastre. Queda prohibido el uso de grasas, aceites o jabon comercial.
  \item \textbf{Empalmes:} Se realizaran unicamente dentro de las cajas de paso u octogonales. Los empalmes se haran mediante uniones firmes soldadas o con conectores a presion, aislados con cinta vulcanizada y protegidos exteriormente con cinta aislante de PVC.
\end{itemize}

\section{CONSTRUCCION DEL POZO DE PUESTA A TIERRA}

El pozo de puesta a tierra vertical se ejecutara de acuerdo a las siguientes especificaciones constructivas:
\begin{enumerate}
  \item \textbf{Excavacion:} Se realizara una excavacion cilindrica o cuadrada de \(0.80\text{ m}\) de diametro y una profundidad minima de \(2.60\text{ m}\).
  \item \textbf{Preparacion del terreno:} El terreno extraido se cernira y mezclara con aditivos quimicos conductivos (dosis de gel o sales especiales) para reducir la resistencia especifica del suelo.
  \item \textbf{Instalacion del electrodo:} Se colocara la varilla de copperweld de \(3/4'' \times 2.40\text{ m}\) de forma perfectamente vertical en el centro del pozo.
  \item \textbf{Relleno y compactacion:} Se rellenara el pozo compactando la tierra tratada por capas de \(0.30\text{ m}\), humedeciendo abundantemente cada capa.
  \item \textbf{Conexion:} El conductor de cobre desnudo de \(10\text{ mm}^2\) se unira a la varilla mediante una abrazadera de bronce tipo perno en U.
  \item \textbf{Caja de registro:} En la parte superior se colocara una caja de concreto prefabricada con tapa, la cual albergara la conexion y permitira realizar las mediciones del protocolo de pruebas.
\end{enumerate}

\section{PRUEBAS DE PUESTA EN SERVICIO}

Una vez terminada la instalacion, el contratista o estudiante realizara las siguientes pruebas en presencia del supervisor:
\begin{itemize}
  \item \textbf{Prueba de continuidad:} Verificacion de la continuidad electrica en todos los circuitos de alumbrado y conductores de puesta a tierra.
  \item \textbf{Medicion de aislamiento:} Utilizando un megometro de 500 V DC, se medira la resistencia de aislamiento entre conductores activos y entre conductor activo y tierra. Los valores deben ser superiores a 1.0 M\(\Omega\).
  \item \textbf{Medicion de la puesta a tierra:} Utilizando un telurimetro, se medira la resistencia del pozo a tierra mediante el metodo de caida de potencial, la cual debe arrojar un valor menor a 15 Ohmios.
  \item \textbf{Prueba de disparo diferencial:} Verificacion del funcionamiento de los interruptores diferenciales mediante el boton de prueba (Test) integrado.
\end{itemize}
